\section{Private decisions and independently validated measurements}
\label{s5:private-logger}
An arrears officer may see information absent from the warehouse. Adding
states to a predictive HMM does not explain that information. A tractable
alternative obtains an independent labelled measurement sample and models
the action mechanism explicitly. Let $S_t$ have $K$ states, $O_t$ have $M$
signal categories, and $A_t$ have two actions. On one decision clock assume
\[
 O_t\mid S_t\sim E_{\cdot S_t},\quad
 A_t\mid S_t\sim q_{S_t},\quad
 R_t\mid S_t,A_t\sim N(r_{A_t,S_t},\sigma_R^2),\quad
 S_{t+1}\mid S_t,A_t\sim T_{A_t}(S_t,\cdot).
\]
The displayed factors assert conditional independence of signal, logger
randomness and reward noise given the indicated parents. Reward occurs before
the next signal. Conditional independence, time invariance, state labels and
the action/transition modularity are scientific premises. The validation
sample observes labelled $(S,O)$ under the same emission law; labels sampled
uniformly across states need not reproduce the population state frequencies.
Their conditional emission likelihood is the relevant contribution.

\subsection{Identification from unnormalized marked moments}
Let $\mu_s=P(S_t=s)$ and define the observable matrix and reward vector
\[
 J_a(o,o')=P(O_t=o,A_t=a,O_{t+1}=o'),\qquad
 k_a(o)=\E[R_t\mathbf1\{O_t=o,A_t=a\}].
\]
The factorization implies
\[
 J_a=E\operatorname{diag}(\mu\odot q(a))T_aE',\qquad
 k_a=E\operatorname{diag}(\mu\odot q(a))r_a.
\]
If the validated $M\times K$ emission matrix has full column rank, let
$L=(E'E)^{-1}E'$. Then
\begin{equation}
 B_a=LJ_aL'=\operatorname{diag}(\mu\odot q(a))T_a,
 \quad w_a=B_a\mathbf1,\quad
 \mu=\sum_a w_a,\quad q_s(a)=w_{as}/\mu_s.
 \label{s5:marked-identification}
\end{equation}
Positive $w_{as}$ permits recovery of $T_a$ by row division and
$r_a=\operatorname{diag}(w_a)^{-1}Lk_a$. This is a constructive sufficient
identification proof using population moments. Conditioning $J_a$ on $A=a$
without retaining $P(A=a)$ discards the action-specific state mass required
by the inversion. An unvalidated emission law may be identified by other
longer-history restrictions; failure of this particular route is not a
universal latent-identification impossibility theorem.

The smallest singular value of $E$ and the minimum state-action mass control
error amplification. Perturbing $J_a$ with $E$ fixed gives
$\|dB_a\|\le\|L\|^2\|dJ_a\|$; perturbing $E$ changes both inverses.
Rank diagnostics can reveal weak separation, but a labelled audit from a
different measurement process may give a precise, wrong $E$. Validation
invariance and label meaning therefore remain explicit restrictions.
More data from a biased validation source do not remove that bias.

\subsection{Joint likelihood, score and two-sample uncertainty}
For one observed trajectory the likelihood is
\[
 L_i(\vartheta)=\sum_{s_{0:H-1}}\mu_{s_0}
 \prod_{t=0}^{H-1}E_{o_t s_t}q_{s_t}(a_t)
       f_R(r_t\mid s_t,a_t)
 \prod_{t=0}^{H-2}T_{a_t}(s_t,s_{t+1}).
\]
The joint log likelihood adds $\sum_{j=1}^{M_v}\log E_{o_j^v s_j^v}$
from the independent validation sample. Softmax row coordinates enforce
probability constraints; Gaussian reward means are scaled by their known
standard deviation for numerical conditioning. Forward probabilities are
normalized at every step. Reward log-density maxima are subtracted and
added back to the log likelihood, so the stabilized computation evaluates
the same function.

Differentiating the finite hidden-path sum gives
$\nabla\log L_i=\E_\vartheta[\nabla\log L_i^{\rm complete}\mid O,A,R]$.
For transition-logit coordinate $(a,s,j)$, for example, the score is
\[
 \sum_t\mathbf1\{A_t=a\}
 \{\xi_t(s,j)-\gamma_t(s)T_a(s,j)\},
\]
where $\gamma,\xi$ are state and adjacent-state smoothing probabilities.
Emission and logger coordinates have the same expected-count minus
expected-total-times-probability form. Reward-mean scores are posterior
weighted residuals divided by $\sigma_R^2$. The validation emission score
uses the observed state label, not a smoothing probability. Exhaustive
hidden-path enumeration checks the likelihood and a separate derivative
checks this analytic score to below $10^{-9}$ in the small reference.

Let trajectory and validation scores be $s_i$ and $u_j$. With fixed horizon,
an identified interior parameter, differentiability, finite moments and
uniform local information, Taylor expansion of the total score gives
\[
 \hat\vartheta-\vartheta_0=
 A^{-1}\left(\sum_{i=1}^Ns_i+\sum_{j=1}^{M_v}u_j\right)
 +o_p(N^{-1/2}),\qquad
 A=N I_o+M_v I_v,
\]
when $M_v/N$ converges to a positive finite constant. Under the correctly
specified model the covariance is $A^{-1}$. A two-sample sandwich uses
$A^{-1}(N J_o+M_v J_v)A^{-1}$, centering each independent score sample
separately. It estimates local sampling uncertainty under possible
likelihood misspecification; it does not restore a causal interpretation
when the model factorization is wrong. If validation alone identifies a
direction, letting $M_v/N\to0$ may change that direction's rate.

Partitioning information into structural and emission parameters shows why
plugging in calibration as known can understate variance. The effective
structural information is the Schur complement
$A_{\beta\beta}-A_{\beta E}A_{EE}^{-1}A_{E\beta}$, not
$A_{\beta\beta}$. For a fixed signal policy $\pi(a\mid o)$,
$\bar\pi(a\mid s)=\sum_oE_{os}\pi(a\mid o)$ determines the target
transition and reward. The value gradient therefore includes the emission
derivative as well as transition, reward and initial-state derivatives.
The historical logger disappears from the intervention recursion but remains
in the likelihood that identifies those parameters.

\subsection{Numerical and inferential challengers}
The new finite experiment fits three states, four signals, private logging,
eight decisions, 600 customer trajectories and either 90 or 600 validation
labels. It estimates 32 parameters using two data-based starts. Reward noise
standard deviation is known at 12. The true fixed observable-policy contrast
is $59.5645$. With 90 labels, bias is $2.0798$ and coverage .850; with 600,
bias is $.1926$ and coverage .875. These denominators include all 80 requested
replications per case, including four and five boundary/numerical failures.
Those failures yield unavailable inference. Model, sandwich and
fixed-emission approximations happen to give the same coverage here; the
reward is informative about state and the measured calibration contribution
is small. This experiment does not demonstrate a large calibration variance
penalty. Omitting the private logger on the first 20 datasets yields biases
of $-36.51$ and $-41.12$, respectively.

The first 600-label dataset was selected before a 199-refit parametric
bootstrap challenger. It resimulates complete trajectories and independent
emission measurements conditional on the audited label counts, then refits
the complete model. Seventeen refits hit a boundary or failed, exceeding the
prespecified 5\% availability limit, so no percentile/basic interval is
released. The standard deviation among 182 usable refits is $10.20$, only a
descriptive conditional diagnostic. Extra starts change the fitted contrast
by $.0012$ and the average negative log likelihood by less than $10^{-10}$.
Information differencing steps $3\times10^{-5},10^{-4},3\times10^{-4}$
change the model standard error by less than $4\times10^{-8}$. The local
information condition number is about 1,929. These checks make numerical
roundoff an implausible sole explanation of the observed undercoverage;
they do not prove the global optimum or finite-sample normality.

The method decision is therefore explicit: this validated finite route is
an identified research fit with a useful value functional, but its current
normal/bootstrapped interval is not a routine calibrated product default.
The larger-sample and audited-trajectory experiments in
Section~\ref{s6:servicing} provide selected alternatives; they do not
retroactively qualify this 600-trajectory interval. Weak-state/profile
inference and bank measurement validation remain distinct requirements. The observed
failure is part of the proposal's method hierarchy, not a hidden exception.

\subsection{Proxies identify a different object with different information}
Sequential proxy bridges can identify value without recovering a uniquely
labelled latent transition model \citep{bennettprox2023,zhangprox2024}.
Their conditional independence and completeness restrictions are different
from labelled measurement validation. The two-stage dynamic-regime source
identifies joint counterfactual quantities along compatible treatment
histories under its stated bridge restrictions; it does not identify every
cross-world joint law. Bridge existence must be asserted or proved separately
from uniqueness. For an operator $K$, injectivity means $Kh=0$ implies
$h=0$; it does not mean that every observed right-hand side lies in
$\operatorname{range}(K)$.

This matters for missing rewards. The 2026 missingness-aware MDP preprint
uses augmented state, a shadow next-state variable, exclusion and completeness
to motivate a bridge \citep[Assumptions 3.1--3.4, Theorem 3.6]{wei2026}.
Its general existence discussion cannot be adopted from injectivity alone.
The operator contract therefore needs an integrable bridge solution and
the relevant range condition before uniqueness can justify identification.
An inverse-problem regularization analysis is then needed to control
estimation error. In a finite square full-rank system
existence follows for every right-hand side; an infinite-dimensional
conditional expectation operator needs more. Overlap, residual and rank
diagnostics are conditional checks. They cannot certify the unobserved
exclusion and completeness assumptions that make an MNAR query causal.
